\section{Weak frequency constraints and the physical space}
\label{sec:physical}
\subsection{The single-observer frequency split}
Let $\Psi_n^s$ be the normalized positive-frequency modes of~\cite{KL},
and $\widetilde\Psi_n^s(t,\theta)=(-1)^n\Psi_n^s(-t,\theta)$.
Appendix~\ref{app:normalization} proves their wave-packet normalization
and mutual orthogonality by radial unitary evolution. They span the full
continuous multiplicity-two spectrum at each angular label. Together with
the bound modes of \eqref{eq:darboux-spectrum}, they form the core of
Theorem~\ref{thm:ambient-core}.

Let $H_+\subset L^2(X)$ be the original positive-frequency subspace and
$H_\perp=H_+^\perp$. Both are $G$-invariant. Write $\cS_+$ for the finite
span of positive-frequency packets with coefficients in $\Cc((0,\infty))$
and $\cS_\perp$ for the negative-frequency packets and bound modes. For
$\mathcal K_N=\Sym^N P_\nu$ define
\begin{equation}
 \cE_{\epsilon\delta}=\operatorname{span}\{u_Lu_Rv:
 u_L\in\cS_\epsilon,\ u_R\in\cS_\delta,\ v\in\mathcal K_N\},
 \quad \epsilon,\delta\in\{+,\perp\}.
\end{equation}
Their $B$ images are mutually orthogonal: between distinct sectors one
observer factor in each group matrix coefficient is zero. Absolute group
averaging and Theorem~\ref{thm:ambient-core} give
\begin{equation}
 \cF(\mathcal K_N)=\cF_{++}\oplus\cF_{+\perp}
 \oplus\cF_{\perp+}\oplus\cF_{\perp\perp},\quad
 \cF_{\epsilon\delta}=\overline{B\cE_{\epsilon\delta}}.
 \label{eq:four-sectors}
\end{equation}
This constructs a bounded physical frequency projection. It does not
evaluate a kinematical projection at a measure-zero trivial representation
in a regular Plancherel decomposition.

\begin{definition}[Weak positive frequency]
A state $k\in\cF(\mathcal K_N)$ is weakly positive frequency in both
observers if
\begin{equation}
 T_k(\phi)=0\quad\text{for every }\phi\in
 \cE_{+\perp}+\cE_{\perp+}+\cE_{\perp\perp}.
 \label{eq:weak-positive}
\end{equation}
The pairings are the absolute position integrals \eqref{eq:position-pairing}.
\end{definition}
This definition uses only single-observer frequency and position tests.
By \eqref{eq:four-sectors}, its solution space is exactly $\cF_{++}$.

\subsection{Actual product amplitudes and all multiplicities}
For $N\ge1$, complete normalized CG analysis gives at fixed actual
right-observer label $a$
\begin{equation}
 (P_a\otimes\Sym^NP_\nu)^{\mathrm P}
 \simeq\int_{d>0}^{\oplus}\dd d\,(P_d\otimes\cM_{N,d,a}).
 \label{eq:cg-multiplicity}
\end{equation}
Inductively add one principal-series particle at a time. Only
principal--principal and principal--discrete decompositions are needed.
The final discrete summands are retained separately; internal discrete
paths remain in $\cM$. Bose symmetrization is a decomposable orthogonal
projection on multiplicity, creating no new spectral type. The original
magnetic basis trivializes the measurable field in $a$; explicit CG
coefficients and their removable limits give jointly measurable transforms.

Each $P_d$ has one unit $K=0$ vector, while a final discrete representation
has none. Hence
\begin{equation}
 U_{N,a}:(P_a\otimes\Sym^NP_\nu)^{(0)}
 \longrightarrow\int_{d>0}^{\oplus}\dd d\,\cM_{N,d,a}
\end{equation}
is unitary. Take a countable original finite-mode Bose basis $e_i(a)$ and
write $u_i(a,d)=U_{N,a}e_i(a)$, so $\int\norm{u_i(a,d)}^2\dd d=1$.

For a compact smooth left packet $f$ of magnetic label zero and an
arbitrary right total-$K=0$ vector $v$, the group-averaged amplitude obeys
\begin{equation}
 g_{f,v}(d)=\KP(d)f(d)v(d),\qquad
 \norm{g_{f,v}}_{\spec}\le B_f\norm v,\qquad
 B_f^2=\esssup_{d>0}\KP(d)|f(d)|^2.
 \label{eq:fixed-left-bound}
\end{equation}
The left mass integration must be performed first. Indeed
\begin{equation}
 |D^s_{00}(g)|\le\Xi_0(\lambda):=P_{-1/2}(\cosh\lambda)
 \le C(1+\lambda)e^{-\lambda/2}
\end{equation}
uniformly in $s>0$, by the circle integral of the Legendre function.
Compact left spectral integration gives the left coefficient $L(g)$
decay $e^{-\lambda/2}(1+\lambda)^{-r}$ up to exponentially smaller terms.
For $r\ge3$, $\int_G|L|\Xi_0\dd g<\infty$. The right cross-spectral
density is in $L^1(\dd d)$ by Cauchy--Schwarz, so Fubini is absolute.
The wave-packet Schur relation gives \eqref{eq:fixed-left-bound};
no absolute integration at two fixed masses is asserted.

For the actual original product $v(a)=k_R(a)e_i(a)$,
\begin{equation}
 g_{i;f,k_R}(d,a)=\KP(d)f(d)k_R(a)u_i(a,d).
 \label{eq:actual-amplitude}
\end{equation}
These amplitudes are dense in
\begin{equation}
 \cH_N^{\spec,0}=\int_{d>0}^{\oplus}\rhoP(d)\dd d
                  \int_{a>0}^{\oplus}\dd a\,\cM_{N,d,a}.
 \label{eq:full-physical-spectrum}
\end{equation}
To prove this, take $h$ orthogonal to all \eqref{eq:actual-amplitude}.
Since $\rhoP\KP=1$, we have
\[
 \int f(d)k_R(a)\ip{h(d,a)}{u_i(a,d)}\dd a\dd d=0.
\]
On compact $A\times D$ the absolute integral is bounded by
$\norm h_{\spec}\sqrt{|A|\sup_D\rhoP^{-1}}$.
Product-test totality implies that the pairing vanishes almost everywhere
for all countably many $i$ simultaneously. Restrict $h$ to $A\times D$
and apply $U_{N,a}^{-1}$; every original basis coefficient vanishes.
Compact exhaustion gives $h=0$, without a bounded diagonal-trace map.

For all finite left magnetic labels the amplitude is
\begin{equation}
 g_\psi(d)=\KP(d)\sum_j(-1)^{m_j}f_j(d)v^{\mathrm P}_{j,-m_j}(d).
 \label{eq:finite-magnetic-amplitude}
\end{equation}
These terms add coherently. In the compact picture every unit
magnetic-basis matrix coefficient is bounded by $\Xi_0$, so the same
left-packet estimate applies. At fixed right total magnetic label only
finitely many final discrete representations occur; exponential decay
makes their mixed integrals absolute, and cross-series Schur terms vanish.
Both norms and cross-pairings of the full product class therefore equal
the spectral inner product. The dense small class
\eqref{eq:actual-amplitude} has the same FP closure as $B\cE_{++}$.

For $N=0$ the right total label is $d=a$. The spectral space is
$L^2(\rhoP(d)\dd d)$ and $g_{f,k_R}=\KP f k_R$ has dense image.
There is no independent $\dd a$ integral. We obtain a concrete
surjective isometry $J_N^0:\cH_N^{\spec,0}\to\cF_{++}$ for every finite
$N$, specified on actual seeds by $J_N^0g_\psi=B\psi$.

\subsection{Observer thresholds and maximality}
Let $\cS_{+,<}$ contain positive-frequency packets supported in
$(0,s_*)$. In addition to \eqref{eq:weak-positive}, require
\begin{equation}
 T_k(u_Lu_Rv)=0
 \label{eq:weak-threshold}
\end{equation}
when either $u_L\in\cS_{+,<}$ and $u_R\in\cS_+$ or conversely.
Define $\cH_N^{\max,w}(s_*)$ to contain all finite FP states satisfying
these two independent test conditions.

\begin{theorem}[Maximal weak-frequency realization]
Identity \eqref{eq:main-space} holds as an equality of concrete subspaces
of $\cF(\mathcal K_N)$. For $N\ge1$, its spectral space is
\eqref{eq:full-physical-spectrum} restricted to $a,d\ge s_*$. For $N=0$
only $d\ge s_*$ remains.
\label{thm:maximal-space}
\end{theorem}
\begin{proof}
Equation \eqref{eq:four-sectors} first puts the state in
$\cF_{++}=J_N^0\cH_N^{\spec,0}$. Use \eqref{eq:actual-amplitude} with
$f$ supported below $s_*$ and repeat the local $L^1$/inverse-CG argument
to obtain $g=0$ on $d<s_*$. Taking instead $k_R$ below $s_*$ gives
$g=0$ on $a<s_*$. Conversely, \eqref{eq:finite-magnetic-amplitude}
retains left support as $d$ support and right support as $a$ support,
since CG analysis is at fixed $a$. Every low-observer-spectrum test is
therefore orthogonal to a high-supported $g$. The fixed-threshold actual
product construction has dense image in exactly that spectral subspace.
The vacuum follows from $\KP f k_R$. Threshold atoms are not additional
vectors in this identified $L^2$ space.
\end{proof}

The topology qualification is substantive. Equation (4.40) of~\cite{KL}
does not prescribe a test topology for an arbitrary non-normalizable
coefficient sum. If it admits the absolute test pairings here with the
stated positive-frequency support, Theorem~\ref{thm:maximal-space}
applies. It is not a theorem about every formal summation prescription.

There is also a sufficient limit criterion. Suppose
$k_j\in\cH_N^{\gen}(s_*)$, $\sup_j\norm{k_j}_{\FP}<\infty$, and
$Rk_j\to F$ locally as distributions. Hilbert weak compactness gives a
weakly convergent subsequence. Lemma~\ref{lem:position-distribution}
identifies its limit as $F=Rk$. The generated space is weakly closed,
so $k$ belongs to it; injectivity makes every weak cluster point the same.
The full sequence converges weakly and
$\norm k\le\liminf_j\norm{k_j}$. Convergence of the norms to $\norm k$
gives strong convergence. A local distributional limit without a uniform
FP bound is not covered.
